\documentclass{article}
\usepackage{mathtools} 
\usepackage{fontspec}
\usepackage[UTF8]{ctex}
\usepackage{amsthm}
\usepackage{mdframed}
\usepackage{xcolor}
\usepackage{amssymb}
\usepackage{amsmath}
\usepackage{hyperref}


% 定义新的带灰色背景的说明环境 zremark
\newmdtheoremenv[
  backgroundcolor=gray!10,
  % 边框与背景一致，边框线会消失
  linecolor=gray!10
]{zremark}{注释}

% 通用矩阵命令: \flexmatrix{矩阵名}{元素符号}{行数}{列数}
\newcommand{\flexmatrix}[4]{
  \[
  #1 = \begin{pmatrix}
    #2_{11}     & #2_{12}     & \cdots & #2_{1#4}   \\
    #2_{21}     & #2_{22}     & \cdots & #2_{2#4}   \\
    \vdots      & \vdots      & \ddots & \vdots     \\
    #2_{#31}    & #2_{#32}    & \cdots & #2_{#3#4}
  \end{pmatrix}
  \]
}

% 简化版命令（默认矩阵名为A，元素符号为a）: \quickmatrix{行数}{列数}
\newcommand{\quickmatrix}[2]{\flexmatrix{A}{a}{#1}{#2}}


\begin{document}
\title{5.2注释}
\author{张志聪}
\maketitle

\begin{zremark}
  实二次型$f,g$的规范型在秩相同的情况下，为什么正惯性指数不同就不能是等价的？
\end{zremark}

\textbf{证明：}

反证法，设$f,g$的正惯性指数不同，但$f, g$等价。设$f, g$的正惯性指数分别为$p, q$且$p \neq q$，
并设它们的矩阵为$A, B$。

因为$f, g$等价，于是存在可逆线性变数替换，使得
\begin{align*}
  B = T^T A T
\end{align*}
于是可得$f, g$分别有两种规范型表示，这与规范型的唯一性矛盾。


\end{document}
